\color{unimain}{\subsection{Защита от неполнофазного режима}}
\color{black}

\begin{enumerate}[label=\arabic{section}.\arabic{subsection}.\arabic*,  align=left, labelsep=4pt, leftmargin=0pt, itemindent=57pt]

\item
Защита предназначена для ликвидации неполнофазного режима работы выключателей с
пофазным приводом, возникающего в результате возможного отказа переключения одной или двух фаз
выключателя при коммутации.

\item
Ввод в работу ЗНР осуществляется с помощью уставки \textbf{«Работа ЗНР»}. Оперативный
ввод/вывод ЗНР производится с помощью БУ (блока управления функциями) \textit{«Вывод ЗНР»}. Пуск ЗНР
формируется при условии появления конфигурируемых сигналов \textit{«ЗНФ В1»} или \textit{«ЗНФ В2»} и
отключенном положении смежного выключателя (для схем с двумя выключателями), с контролем
протекания тока нулевой или обратной последовательности. Выбор количества контролируемых
выключателей задается уставкой \textbf{«ЗНР контр. В2»}, \textbf{«ЗНР контр. В3»}. Порог срабатывания ЗНР по
току нулевой последовательности задается уставкой \textbf{«Iср ЗНР»}. Выдержка времени срабатывания ЗНР
задается уставкой \textbf{«Тср ЗНР»}. 

\begin{figure}[H]
\centering
\includegraphics[width=0.7\textwidth,height=0.7\textheight,keepaspectratio]{\fbpath/161.0221 ЗНР/img/ZNR.pdf}
\captionof{figure}{Функциональный блок ЗНР}\label{pict_znr:ZNR}
\end{figure}

\small
\begin{longtable}{|>{\centering\arraybackslash}m{5.3cm}|>{\centering\arraybackslash}m{4.2cm}|>{\centering\arraybackslash}m{1cm}|>{\centering\arraybackslash}m{3.3cm}|>{\centering\arraybackslash}m{1.8cm}|}
% Первая страница: полный заголовок
\caption{Перечень уставочных параметров ЗНР\hfill\vspace{-0.5\baselineskip}}\label{znr:tbl1}\\
\hline
\rowcolor{gray!30}
Наименование уставки & Диапазон & Шаг & Значение по умолчанию & Единица измерения \\
\hline
\endfirsthead

% Промежуточные страницы: «Продолжение таблицы»
\caption*{\hspace{3pt}\emph{Продолжение таблицы \ref{znr:tbl1}}\hfill\ }\vspace{-0.5\baselineskip} \\
\hline
\rowcolor{gray!30}
Наименование уставки & Диапазон & Шаг & Значение по умолчанию & Единица измерения \\
\hline
\endhead

%===m>ЗНР|-
%===m
\end{longtable}
\normalsize

\end{enumerate}